% HAND-WRITTEN fragment -- NOT generated by maestro2tex, and Adc.tex is itself
% hand-written; kept separate so the section frame stays untouched. Input from
% AnalogChapter.tex immediately after Adc.tex.
%
% Salvaged from the retired note_calibration_seq.tex step 3 (converter zero) and
% note_calibration_residual.tex (the dither argument). Fabric from
% ~/vestarv/docs/afe_rev2_debug_mux.html §6.3: AFE_ADCSRCSEL[2:0], 0 = vout,
% 1 = vcm, 2 = vref, 3 = RE, 4 = ATB; §9 step 10 is the external ramp;
% per-channel slot 9 = adc_in, observable between conversions only.
% Numbers: tab_adc_transfer (1.535x slope, 0.441 V intercept, 1.591 mV LSB);
% DAC step 528.96 uV (tab_dacr2r12_mc_reference). No cell names in prose.

\subsubsection{Measuring this block on chip}

The converter is the one block whose gain has never been checked in system, and
the one that can be characterised end to end with no analog pad measurement.
Set \texttt{AFE\_ADCSRCSEL}~=~4 to route the test bus into the converter input,
drive an external ramp onto \texttt{atb\_d} with \texttt{AFE\_DBGDIREN}~=~1 and
\texttt{AFE\_DBGBUFEN}~=~0, and fit slope and intercept per die: the
\(1.535\times\) slope, \SI{0.441}{\volt} code-zero intercept and
\SI{1.591}{\milli\volt} effective LSB of Table~\ref{tab:adc-transfer} rest on a
single nominal extracted run with no chip-to-chip data behind them. For the zero
code, set \texttt{AFE\_ADCSRCSEL}~=~1 and average the common-mode rail; that
code is the origin of the current axis and carries the converter's offset and
the common-mode value together, which is the combination a current reading
needs. Averaging reaches the quantisation floor only if the input dithers, and
with the electrode isolated the amplifier output noise is tens of microvolts ---
so step the common-mode DAC over three of its \SI{529}{\micro\volt} codes, which
span one converter code. Observe the sampling node itself at per-channel slot 9
between conversions only. This test fabric is design intent; it is not yet in
the schematics.
